
\documentclass{article}
\usepackage{colortbl}
\usepackage{makecell}
\usepackage{multirow}
\usepackage{supertabular}

\begin{document}

\newcounter{utterance}

\centering \large Interaction Transcript for game `cladder', experiment `full\_v1.5\_default', episode 4591 with Qwen3.6{-}27B.
\vspace{24pt}

{ \footnotesize  \setcounter{utterance}{1}
\setlength{\tabcolsep}{0pt}
\begin{supertabular}{c@{$\;$}|p{.15\linewidth}@{}p{.15\linewidth}p{.15\linewidth}p{.15\linewidth}p{.15\linewidth}p{.15\linewidth}}
    \# & \multicolumn{2}{c}{Player} && \multicolumn{2}{c}{Game Master} \\
    \hline

    \theutterance \stepcounter{utterance}  
    & & & \multicolumn{4}{p{0.6\linewidth}}{
        \cellcolor[rgb]{0.9,0.9,0.9}{
            \makecell[{{p{\linewidth}}}]{
                \texttt{\tiny{[P1$\langle$GM]}}
                \texttt{You are an expert in causal inference. The following question is not a typical commonsense query, but rather a meticulously designed question created by a professor specializing in causal inference, intended to assess the students' mastery of the course content.} \\
\texttt{Imagine a self{-}contained, hypothetical world with only the following conditions, and without any unmentioned factors or causal relationships: Having a brother has a direct effect on wife and alarm clock. Wife has a direct effect on alarm clock. The overall probability of having a brother is 1\%. The probability of not having a brother and ringing alarm is 64\%. The probability of having a brother and ringing alarm is 0\%. Is the chance of ringing alarm larger when observing having a brother?} \\
\texttt{Start your answer with 'Yes' or 'No', followed by additional reasoning or evidence to support your explanation.} \\
            }
        }
    }
     \\ \\

    \theutterance \stepcounter{utterance}  
    & \multicolumn{4}{p{0.6\linewidth}}{
        \cellcolor[rgb]{0.561,0.737,0.561}{
            \makecell[{{p{\linewidth}}}]{
                \texttt{\tiny{[P1$\rangle$GM]}}
                \texttt{No} \\
\\ 
\texttt{**Reasoning:**} \\
\\ 
\texttt{To determine whether the chance of ringing the alarm is larger when observing that one has a brother, we need to compare the conditional probability of the alarm ringing given that one has a brother, \$P(\textbackslash{}text\{Alarm\} | \textbackslash{}text\{Brother\})\$, with the conditional probability of the alarm ringing given that one does not have a brother, \$P(\textbackslash{}text\{Alarm\} | \textbackslash{}neg \textbackslash{}text\{Brother\})\$. Alternatively, we can simply compare these conditional probabilities to the marginal probability or to each other directly.} \\
\\ 
\texttt{Let \$B\$ denote the event "Having a brother" and \$A\$ denote the event "Ringing alarm".} \\
\\ 
\texttt{From the problem description, we are given:} \\
\texttt{1.  \$P(B) = 0.01\$ (1\%)} \\
\texttt{2.  Therefore, \$P(\textbackslash{}neg B) = 1 {-} P(B) = 0.99\$} \\
\texttt{3.  \$P(\textbackslash{}neg B \textbackslash{}cap A) = 0.64\$ (The probability of not having a brother and ringing alarm is 64\%)} \\
\texttt{4.  \$P(B \textbackslash{}cap A) = 0\$ (The probability of having a brother and ringing alarm is 0\%)} \\
\\ 
\texttt{We calculate the conditional probabilities:} \\
\\ 
\texttt{**1. Probability of alarm ringing given having a brother:**} \\
\texttt{\$\$P(A | B) = \textbackslash{}frac\{P(B \textbackslash{}cap A)\}\{P(B)\} = \textbackslash{}frac} \\
            }
        }
    }
    & & \\ \\

    \theutterance \stepcounter{utterance}  
    & & & \multicolumn{2}{p{0.3\linewidth}}{
        \cellcolor[rgb]{0.9,0.9,0.9}{
            \makecell[{{p{\linewidth}}}]{
                \texttt{\tiny{[GM$|$GM]}}
                \texttt{No} \\
            }
        }
    }
    & & \\ \\

    \theutterance \stepcounter{utterance}  
    & & & \multicolumn{2}{p{0.3\linewidth}}{
        \cellcolor[rgb]{0.9,0.9,0.9}{
            \makecell[{{p{\linewidth}}}]{
                \texttt{\tiny{[GM$|$GM]}}
                \texttt{no} \\
            }
        }
    }
    & & \\ \\

    \theutterance \stepcounter{utterance}  
    & & & \multicolumn{2}{p{0.3\linewidth}}{
        \cellcolor[rgb]{0.9,0.9,0.9}{
            \makecell[{{p{\linewidth}}}]{
                \texttt{\tiny{[GM$|$GM]}}
                \texttt{game\_result = WIN} \\
            }
        }
    }
    & & \\ \\

\end{supertabular}
}

\end{document}
